% \iffalse meta-comment
%
% hit-report-matter.dtx 是报告的文档骨架
%
% \fi
%
% \section{报告的前置、正文与后置}
%
% \changes{v4.0a}{2026/09/28}{报告也有 \cs{frontmatter}、\cs{mainmatter}、\cs{backmatter}，只切页码与版式，
%   由第一个前置部件（目录）、正文第一个标题、第一个后置部件（参考文献）自动切换；深圳博士中期的边框
%   与深圳硕士的页眉在进正文时挂上；学位论文专有的部件在报告里写了也认，判断不排，不再报未定义}
% \changes{v4.0a}{2026/09/14}{文献表调总出口 \cs{hit@bib@print:}；正文照终稿排的那一档页码制式跟着终稿；
%   顶级另起一页只给深圳本科}
%
% ^^A ======================================================================
% ^^A Module report-matter: 报告的前后文切换（hit-report）
% ^^A ======================================================================
%
% 与学位论文同名的三个命令，写法一样；部件排什么由部件自己管，这三个只切页码与版式。
% 报告的前置只有封面与目录，后置只有参考文献。
%
%    \begin{macrocode}
%<*report-matter>
\bool_if:NF \g__hit_final_bool
  {
%    \end{macrocode}
%
% \emph{页码制式跟着正文那一档走。} 报告默认整篇一个制式；正文照终稿排的那一档
% 与终稿一样分两段——前置（目录）罗马数字，正文另起一段从 1 数。原件的页脚
% 逐页量过就是这样：目录印 \texttt{I}、正文印 \texttt{-1-}。
%
% 章另起一页也在这儿开：这一档正文的顶级按章排（见 \file{src/hithesis.dtx} 的层级
% 错位），章就得另起一页。顶级可能是 \cs{section} 也可能是 \cs{chapter}（见
% \file{src/flow/hit-report-heading.dtx} 的标题层级），所以不往哪个命令上挂补丁，
% 只立 \cs{g\_\_hit\_report\_top\_break\_bool}，由顶级那一支换页。
% 本部与威海报告的顶层节标题不另起一页，仅深圳本科设置该标志。威海硕士开题
% 照 \issue[219] 的模板，节套的是终稿的“标题 2”，也不另起一页。
% 类自己排的封面与目录按原级走，排的时候立 \cs{g\_\_hit\_report\_heading\_raw\_bool}。
%
% \emph{不能走 \pkg{ctex} 的 \option{break} 键。} 那个键在 \cs{@startsection} 里
% 挂在 \cs{if@nobreak} 的假支上——标题底下没正文时这个开关还立着（\hologo{LaTeX}
% 不在标题与它的头一段之间断页），换页就被跳过。示例里 1.2 底下正好没字，
% 章 2 于是跟在同一页上。所以在顶级标题前面无条件发一次换页。
%
% \emph{换页之后那一段的段前也要留住。} \hologo{TeX} 把换页处的竖胶丢掉，而
% Word 在页首照进段前：原件的章标题基线量到 137.98 ＝ 版心顶 107.75 加段前 12
% 加上升部 18.23，丢了就只剩 125.98。这一档的标题由 \pkg{docxlike} 排，换完页立起它的
% “新一节开头”标记，标题自己留住段前，与终稿的章一样（见
% \file{src/flow/hit-report-heading.dtx}）。\cs{topskip} 归零是当初换页后垫 \cs{hrule} 那一招
% 留下的，可去掉之后章标题在页首的段前就不按“段前减上一段段后”算了（第一章落到 137.9，
% Windows Word 的导出件是 128.2），所以留着。
%
% 后置部分及文末须恢复 topskip，否则 clearpage 保留的零高空盒可能产生空白页。
%
% \emph{别的档整篇一个制式，正文从 1 数。} 原先首页与目录在同一个 \env{titlepage} 环境里，单面文档在
% 环境结束时把页码归 1。首页与目录拆成两个部件以后，归 1 挪到这里；双面的不归，与原先一样。
%
% 深圳研究生报告的目录与正文分别从阿拉伯数字 1 起；中期目录使用无页眉样式。
%    \begin{macrocode}
\dim_new:N \g__hit_report_topskip_dim
\cs_set_protected:Npn \frontmatter
  {
    \hit@matter@set:n { front }
    \bool_if:NT \g__hit_final_body_bool
      {
        \legacy_if_gset_false:n { @mainmatter }
        \pagenumbering
          {
            \bool_lazy_and:nnTF
              { \bool_if_p:N \g__hit_shenzhen_bool }
              { \bool_lazy_or_p:nn{\bool_if_p:N\g__hit_proposal_bool}{\bool_if_p:N\g__hit_shenzhen_report_bool} }
              { arabic } { Roman }
          }
        \pagestyle
          { \bool_if:NTF \g__hit_harbin_report_bool { hit@empty } { hit@headings@front } }
        \bool_if:NT\g__hit_shenzhen_report_bool{\pagestyle{hit@shenzhen@toc}}
      }
  }
\cs_set_protected:Npn \mainmatter
  {
    \hit@matter@set:n { main }
    \bool_if:NTF \g__hit_final_body_bool
      {
        \clearpage
        \legacy_if_gset_true:n { @mainmatter }
        \pagenumbering { arabic }
        \pagestyle { hit@headings }
        \dim_gset_eq:NN \g__hit_report_topskip_dim \topskip
        \dim_zero:N \topskip
        \bool_if:NT\g__hit_shenzhen_report_bool
          {\bool_gset_true:N\g_docxlike_section_start_bool}
        \bool_lazy_and:nnT{\bool_if_p:N\g__hit_shenzhen_bool}{\bool_if_p:N\g__hit_bachelor_bool}
          { \bool_gset_true:N \g__hit_report_top_break_bool }
      }
      { \legacy_if:nF { @twoside } { \setcounter { page } { 1 } } }
  }
\cs_new_protected:Npn \__hit_report_topskip_restore:
  {
    \bool_if:NT \g__hit_final_body_bool
      { \dim_set_eq:NN \topskip \g__hit_report_topskip_dim }
  }
\cs_set_protected:Npn \backmatter
  {
    \hit@matter@set:n { back }
    \__hit_report_topskip_restore:
  }
\AtEndDocument { \hit@matter@if:nT { main } { \__hit_report_topskip_restore: } }
%    \end{macrocode}
%
% 报告识别终稿专有部件并按支持情况跳过。扉页、摘要、
% 符号表、缩略语表、符号与缩略语合页、答辩页、声明、索引在报告里没有版式，\option{shown=true} 也不排，给一条警告；
% 其余的键不认也不报，写给学位论文的参数原样留着就行。环境体整个跳过。\cs{maketitle} 是
% \hologo{LaTeX} 内核的，报告里改成同一个判断；\cs{printindex} 等导言区装完再看，用户自己装了
% 索引宏包的照它的。结论、成果、致谢、简历、永久通讯地址报告里有版式，见
% \file{src/pages/hit-final-backmatter.dtx}。
%    \begin{macrocode}
\keys_define:nn { hit / part / report }
  {
    shown .code:n = { \hit@part@set@shown:n {#1} } ,
    unknown .code:n = { } ,
  }
\cs_new_protected:Npn \__hit_report_final_only:nn #1#2
  {
    \tl_set:Nn \l_hit_part_shown_tl { auto }
    \keys_set:nn { hit / part / report } {#2}
    \str_if_eq:VnT \l_hit_part_shown_tl { true }
      { \msg_warning:nnn { hithesis } { part-final-only } {#1} }
  }
\cs_new_protected:Npn \__hit_titlepage_command:n #1
  { \__hit_report_final_only:nn { titlepage } {#1} }
\NewDocumentCommand \listofabbreviations { s O{} }
  { \__hit_report_final_only:nn { list-of-abbreviations } {#2} }
\NewDocumentCommand \listofsymbols { O{} }
  { \__hit_report_final_only:nn { list-of-symbols } {#1} }
\NewDocumentCommand \printnomenclature { O{} }
  { \__hit_report_final_only:nn { nomenclature } {#1} }
\NewDocumentCommand \defense { +O{} }
  { \__hit_report_final_only:nn { defense } {#1} }
\NewDocumentCommand \declarations { O{} }
  { \__hit_report_final_only:nn { declarations } {#1} }
\cs_new_protected:Npn \__hit_report_final_only_env:nn #1#2
  {
    \NewDocumentEnvironment {#1} { O{} }
      {
        \__hit_report_final_only:nn {#2} {##1}
        \hit@part@skip:
      }
      { }
  }
\clist_map_inline:nn
  {
    { abstract } { abstract } ,
    { symbols } { list-of-symbols } , { symbols* } { list-of-symbols } ,
    { denotation } { list-of-symbols } , { denotation* } { list-of-symbols }
  }
  { \__hit_report_final_only_env:nn #1 }
\AddToHook { begindocument / end }
  {
    \cs_if_exist:NF \printindex
      {
        \NewDocumentCommand \printindex { O{} }
          { \__hit_report_final_only:nn { index } {##1} }
      }
  }
  }
%</report-matter>
%    \end{macrocode}
%
% \Finale
